\chapter{Testing, analysis, debugging, and proof}

\section{Testing as evidence}
A test is an experiment with a programmed oracle. It does not prove that all
possible behaviour is correct. It provides evidence about a stated input space,
environment, and observation. The strongest test suite is not the one with the
largest line count; it is the one that catches important failures and explains
them quickly.

\noindent\begin{tabularx}{\textwidth}{@{}p{0.22\textwidth}p{0.25\textwidth}X@{}}
\toprule
Level & Best at finding & Main cost \\
Unit & Local rule and edge-case defects & Can miss wiring, configuration, and real dependencies \\
Component & Contract and integration defects & More setup and slower diagnostics \\
End to end & Broken user journeys and deployment assumptions & Slow, brittle, and often hard to localize \\
Property or model & Broad invariant violations and state-machine bugs & Requires a good model and failure shrinking \\
\bottomrule
\end{tabularx}

The levels are not a pyramid law. A small system with a strong component test may
need few end-to-end tests. A safety-critical subsystem may need formal analysis.
Choose the mix from failure cost, change rate, and observability.

\section{Test design}
For each behaviour, identify:

\begin{enumerate}
  \item a representative valid input;
  \item boundary values and malformed input;
  \item a state transition or dependency failure;
  \item the observable result and any invariant that must remain true.
\end{enumerate}

Avoid asserting incidental representation. If a client only needs an order ID
and status, do not make it depend on database column order or an internal class.
Tests are also consumers of an interface; over-specified tests make refactoring
expensive without increasing confidence.

\section{Static analysis and types}
Static analysis checks source without executing every path. A type checker can
catch incompatible values; a linter can catch suspicious constructs; a dataflow
tool can identify possible taint or resource leaks. False positives impose cost,
so tune rules and provide an explicit escape hatch with a reason. A passing
static check is evidence, not a proof of correctness.

Types are especially useful when they encode meaningful distinctions. A string
for a customer ID and a string for a password are technically interchangeable but
not semantically interchangeable. A wrapper type or validation boundary can
reduce accidental mixing. Do not create types that obscure a simple domain rule.

\section{Debugging as hypothesis testing}
Debugging is a search through competing explanations. Start with a reproducible
failure and preserve the original evidence. Write the expected and actual
behaviour, then reduce the input or environment while keeping the failure.

A practical loop is:

\begin{enumerate}
  \item observe: capture logs, traces, inputs, versions, and timing;
  \item hypothesize: name a mechanism that could explain the observation;
  \item predict: identify what else should be true if the hypothesis is correct;
  \item test: run the smallest discriminating experiment;
  \item change: fix the cause or update the model, then add a regression check.
\end{enumerate}

Changing several things at once destroys information. The exception is an active
incident where stabilization is more important than diagnosis; record temporary
changes so they can be revisited.

\section{Formal methods in proportion}
Formal methods use mathematical models to specify or analyze systems. They range
from assertions and contracts to model checking, theorem proving, and formally
verified compilation. The benefit is strongest when a small, critical state space
can be modeled precisely. The cost is building and maintaining the model and
connecting it to implementation.

An invariant can be simple:

\[
  \forall i \in \text{order items}: i.quantity > 0 \land i.unit\_price\_cents > 0.
\]

That statement can be enforced by application validation, a database constraint,
and tests. More advanced tools can explore concurrent transitions. Use formal
analysis where it reduces a material risk, not as decoration.

\section{Mutation and negative evidence}
Mutation testing changes code in small ways and asks whether the tests fail. A
surviving mutation points to a missing assertion, but mutation scores are not a
universal quality metric. Similarly, a security test that finds no vulnerability
does not establish security; it establishes only that the tested attack and
observation did not reveal one.

\begin{exercise}
Create a test matrix for an order endpoint with dimensions of authentication,
payload validity, duplicate key, database failure, and client timeout. Choose the
smallest set of tests that covers each important interaction and explain what is
left untested.
\end{exercise}

\section{Chapter summary}
Tests provide scoped evidence. Combine unit, component, journey, property, and
static checks according to risk. Debug by preserving evidence and testing
hypotheses. Use types, assertions, and formal models when they clarify and reduce
important risk rather than because a technique is fashionable.
